\documentclass[11pt]{article}
\usepackage[margin=1in]{geometry}
\usepackage[T1]{fontenc}
\usepackage[utf8]{inputenc}
\usepackage{lmodern,microtype,amsmath,amssymb,amsthm,mathtools,booktabs}
\usepackage[colorlinks=true,linkcolor=blue,urlcolor=blue,citecolor=blue]{hyperref}
\newtheorem{theorem}{Theorem}
\newtheorem{lemma}[theorem]{Lemma}
\newtheorem{proposition}[theorem]{Proposition}
\newtheorem{corollary}[theorem]{Corollary}
\theoremstyle{remark}\newtheorem{remark}[theorem]{Remark}
\newcommand{\R}{\mathbb R}\newcommand{\E}{\mathbb E}\newcommand{\Prb}{\mathbb P}
\newcommand{\relu}{\operatorname{ReLU}}\newcommand{\Var}{\operatorname{Var}}
\newcommand{\supp}{\operatorname{supp}}\newcommand{\dist}{\operatorname{dist}}
\newcommand{\one}{\mathbf1}\newcommand{\F}{\mathbb F}
\DeclareMathOperator{\spanop}{span}
\title{Superpolynomial online mistakes for isotropic\protect\\ Bayesian ReLU and softplus networks}
\author{ChatGPT-6 Astra \and Kaarel H\"anni \and Dmitry Vaintrob\thanks{Authors are listed in alphabetical order.}}
\date{November 2026}
\hypersetup{pdftitle={Superpolynomial online mistakes for isotropic Bayesian ReLU and softplus networks},pdfauthor={ChatGPT-6 Astra; Kaarel Hanni; Dmitry Vaintrob}}
\begin{document}\maketitle
\begin{abstract}
We construct a fixed order of distinct Boolean inputs on which exact
Bayesian prediction with an ordinary one-hidden-layer ReLU network makes
superpolynomially many mistakes in the strong-prior limit. The weights
have independent Gaussian coordinates, with identical variances within
each layer. Before observing labels, the prior is conditioned on an
absolute output margin at every input; training then requires exact fit.
The construction works at every sufficiently large finite hidden width.
For a parameter $m$, the input dimension and the target's Boolean circuit
size are $O(m^2)$. A shallow ReLU teacher has width $O(m^3)$ and
coefficients bounded by that width, while $2^{m/2-o(m)}$ designated
predictions have correct-label posterior probability tending to zero.
Thus the mistake bound is $\exp(\Omega(C^{1/3}))$ in the displayed
teacher width $C$. The same bound holds for standard softplus, with a
teacher of comparable size and an existential learner-width threshold.
The width threshold and calibration prefix are large, and the input
order is deliberately unrepresentative of the uniform distribution.
The theorem concerns cumulative online mistakes, not iid sample complexity.

The proof determines an exact weighted Fourier norm on an unequal-rail
parity code, proves polynomial stability, and uses calibration to exclude
all competitors outside the code. A path-based linear subcode allows
linearly sized rails. A transfer-matrix argument, together with
independence of squarefree radicals, supplies exponentially many distinct
completion costs. Observability outside a suitable affine set and a
linear rounding estimate identify the unique cheapest completion after
each history. The resulting energy gaps apply to the full Gaussian
posterior, including all parameterizations of each function.
\end{abstract}

\section{Model and main result}

A target can have a small circuit and still be learned poorly along an
unfavorable sequence of inputs. The issue for a Bayesian neural network
is whether its actual posterior can repeatedly prefer a wrong completion,
even when the learner is arbitrarily wider than a small teacher. We prove
that this can occur with ordinary ReLU activation and isotropic Gaussian
coordinates. No rule for selecting among tied minimum-energy classifiers
is used.
Section~\ref{sec:softplus} extends the construction to the smooth
activation $\log(1+e^t)$ under the same Gaussian prior and likelihood.

The input space is $X=\{-1,1\}^d$. Write $\bar x=(1,x)$ and consider
\begin{equation}\label{eq:network}
 F_\theta(x)=\sum_{j=1}^N v_j\relu(u_j\cdot\bar x),
 \qquad \relu(t)=\max(t,0),\quad u_j\in\R^{d+1}.
\end{equation}
The first coordinate of $u_j$ is the weight of the constant input. There
are no additional biases. Independently,
\begin{equation}\label{eq:prior}
 (u_j)_i\sim\mathcal N(0,N\sigma^2/d),\qquad
 v_j\sim\mathcal N(0,\sigma^2/N).
\end{equation}
Its density is proportional to $\exp[-K_N(\theta)/(2\sigma^2)]$, where
\begin{equation}\label{eq:energy}
 K_N(\theta)=\frac dN\sum_j\|u_j\|_2^2+N\sum_jv_j^2.
\end{equation}
The prior is first conditioned on $|F_\theta(x)|>1/2$ at every $x\in X$.
After each prediction it is conditioned on the exact revealed label.
Thus a history $H\subseteq X$ with target $f:X\to\{-1,1\}$ conditions
the Gaussian law on
\begin{equation}\label{eq:posterior-cell}
 |F_\theta(x)|>1/2\quad(x\in X),\qquad
 f(x)F_\theta(x)>1/2\quad(x\in H).
\end{equation}
The strong-prior limit is $\sigma\downarrow0$ at each fixed $d,N$.

For the ReLU model, the conclusion also holds for any two fixed positive layer variances
$\alpha\sigma^2$ and $\beta\sigma^2$ in place of \eqref{eq:prior}.
Positive homogeneity gives the lower bound
$2V(F)/\sqrt{\alpha\beta}$, where $V$ is the variation norm defined below.
Balancing any finite representation of cost $B$ attains energy
$2B/\sqrt{\alpha\beta}$. The constructed minimizers attain this bound
at the widths used below, so their strict energy order is unchanged.
In particular, writing $t=\sqrt N\,\sigma$ turns \eqref{eq:prior} exactly into first-layer variance $t^2/d$ and readout variance $t^2/N^2$. At each fixed width the limits $t\downarrow0$ and $\sigma\downarrow0$ are the same.

\begin{theorem}[A fixed sequence with many confident errors]\label{thm:main}
For every sufficiently large integer $m$, put
\[
 \ell=2m-1,\qquad R=2^{42}\ell,\qquad n_0=\ell R,\qquad
 n_{\max}=n_0+6m-2,
\]
\[
 d=2n_0+6m-2,\qquad N_0=2^{m+n_{\max}+1},\qquad
 k_m=\max\left\{2,\left\lceil
 \frac{m+250+7\log_2\ell}{39+\log_2\ell}\right\rceil\right\},
 \qquad q_m=2k_m.
\]
There are a Boolean target $f_m$ on $d$ bits, a fixed permutation of its
input cube, and
\[
 T_m=\left\lfloor
 \frac{2^{m/2}-\sum_{j=0}^{q_m-1}\binom mj}{4\binom m{q_m}}
 \right\rfloor
 =2^{m/2-o(m)}
\]
designated rounds such that, for every finite $N\ge N_0$, the posterior
probability of the correct label at each designated round tends to zero
as $\sigma\downarrow0$.

The target has a bounded-fan-in Boolean circuit with $O(d)=O(m^2)$ gates
and a one-hidden-layer ReLU implementation of width at most
\[
 C_m=(m+1)(n_0+1)+8m-4=\Theta(m^3),
\]
output margin one, and every coefficient bounded by $C_m$. The same input
permutation works for every $N\ge N_0$.

For every $\eta>0$ and fixed such $m,N$, all sufficiently small positive
$\sigma$ therefore have these properties:
\begin{enumerate}
\item Fresh posterior sampling makes at least $T_m$ mistakes with
probability at least $1-\eta$.
\item Posterior-label majority makes a mistake at every designated round.
\item A prediction rule requiring correct-label posterior probability
strictly greater than $2/3$ fails at every designated round.
\end{enumerate}
\end{theorem}

The width threshold is an explicit sufficient bound. Here
$d=\Theta(m^2)$, $C_m=\Theta(m^3)$, and $q_m=O(m/\log m)$, so the
theorem gives $\exp(\Omega(C_m^{1/3}))$ cumulative mistakes.
The width quantifier covers every fixed $N\ge N_0$; it does not claim
the result for every width merely exceeding the teacher width.
The calibration prefix reveals all inputs outside a small union of
logical fibers. This is compatible with good iid learning guarantees.

\section{Variation norm and parity codes}

All Fourier coefficients and $L^2$ norms use the uniform physical cube.
For $A\subseteq[d]$, write $\chi_A(x)=\prod_{i\in A}x_i$ and
$\widehat F(A)=\E[\chi_A F]$. Define
\[
 V(F)=\inf\left\{\sum_j|c_j|\|a_j\|_2:
 F(x)=\sum_jc_j\relu(a_j\cdot\bar x)\text{ on }X\right\}.
\]
The same arguments apply to the cost of a specified finite
representation, without taking this infimum. The arithmetic--geometric
mean inequality gives
\begin{equation}\label{eq:energy-variation}
 K_N(\theta)\ge2\sqrt d\sum_j|v_j|\|u_j\|_2\ge2\sqrt d\,V(F_\theta).
\end{equation}
Conversely any finite representation can be balanced neuron by neuron
so that its energy equals twice $\sqrt d$ times its variation cost.
Unused neurons may be set to zero.

For positive even $n$, define
\begin{equation}\label{eq:parity-cost}
 \kappa_n=\frac{\binom{n-2}{(n-2)/2}}{2^{n-1}},\qquad
 a(n)=\frac{\sqrt n}{\kappa_n}.
\end{equation}
The alternating binomial identity
\[
 2^{-n}\sum_{k=0}^n(-1)^k\binom nk\max(n-2k,0)
 =(-1)^{n/2+1}\kappa_n
\]
shows that a balanced unit ridge on $n$ coordinates attains absolute
parity correlation $\kappa_n/\sqrt n$. Signed averaging over all
input-coordinate flips represents $\chi_A$ with variation cost
$a(|A|)$, using at most $2^{|A|}$ neurons.

\subsection{The full rail-code atom theorem}

For this theorem only, let $p\ge2$ denote a number of disjoint rail
pairs $(B_i^0,B_i^1)$. Their full transversals are
$\widetilde A_t=\bigcup_iB_i^{t_i}$, indexed by $t\in\F_2^p$.
Let $\mathcal D_{\rm full}$ be the span of characters contained in
at least one such transversal, and let $\mathcal A_{\rm full}$ be
the union of their balanced parameter rays, including zero and all
coordinate sign choices. Each nonzero ray has zero bias, uses one
whole rail per pair, and has equal absolute coordinate weights.

\begin{theorem}[Exact unequal-rail norm and stability]\label{thm:rail}
Suppose all rail sizes are even and belong to $[R,R+\Delta]$, where
$R,\Delta$ are positive even integers, $\Delta\ge2$, and
\begin{equation}\label{eq:large-R}
 R\ge2^{40}p\Delta.
\end{equation}
For every affine ReLU atom $h_u(x)=\relu(u\cdot\bar x)$,
\begin{equation}\label{eq:atom-norm}
 Q_{\rm full}(h_u):=\sum_ta(|\widetilde A_t|)
          |\widehat h_u(\widetilde A_t)|\le\|u\|_2.
\end{equation}
Nonzero equality holds exactly on $\mathcal A_{\rm full}$.
Consequently
\begin{equation}\label{eq:exact-code-norm}
 V\left(\sum_tc_t\chi_{\widetilde A_t}\right)
 =\sum_ta(|\widetilde A_t|)|c_t|.
\end{equation}
The parameter and function stability estimates are
\begin{align}
 \dist(u,\mathcal A_{\rm full})^2
 &\le2^{16}R^3\|u\|_2(\|u\|_2-Q_{\rm full}(h_u)),
 \label{eq:parameter-stability}\\
 \|h_u-P_{\mathcal D_{\rm full}}h_u\|_2^2
 &\le2^{16}R^3\|u\|_2(\|u\|_2-Q_{\rm full}(h_u)).
 \label{eq:atom-stability}
\end{align}
\end{theorem}

A complete proof appears in Appendix~\ref{app:rail}. The exact constant
in the norm and the polynomial stability bound will both be used.

\section{A path-based subcode with small physical dimension}
\label{sec:path-code}

Let $m\ge2$ and $\ell=2m-1$. The rail pairs are indexed by the columns
\[
 a_j\in\{e_1,\ldots,e_m\}
       \cup\{e_1+e_2,e_2+e_3,\ldots,e_{m-1}+e_m\}
       \subseteq\F_2^m.
\]
The first $m$ columns are called loops and the others edges. Every
short rail has size $R=2^{42}\ell$. The long rail has size $R+4$ at
a loop and $R+2$ at an edge. Thus Theorem~\ref{thm:rail} applies with
$p=\ell$ and $\Delta=4$.

Retain only the transversals
\[
 A_s=\bigcup_{j=1}^{\ell}B_j^{a_j\cdot s}
       \quad(s\in G:=\F_2^m),\qquad b=\chi_{A_0}.
\]
If $Ds=(s_i+s_{i+1})_{i=1}^{m-1}$ is the path difference map, then
\begin{equation}\label{eq:path-degrees}
 n_s=|A_s|=n_0+4|s|+2|Ds|,\qquad n_0=\ell R,
 \qquad n_s\le n_{\max}=n_0+6m-2.
\end{equation}
The total physical dimension is $d=2n_0+6m-2$.
Let $\xi_j\in\F_2$ encode the parity of the whole $j$th rail pair,
and define the logical syndrome
\[
 (-1)^{\xi_j}=\chi_{B_j^0\cup B_j^1}(x),\qquad
 z=\sum_ja_j\xi_j.
\]
The vector $z$ is uniform because the loop columns span $G$, and $b$
is independent of $z$. Moreover $\chi_{A_s}=b(-1)^{s\cdot z}$.
Let $\mathcal C$ be the span of these retained characters and
$\mathcal D$ their downward Fourier closure. In particular, an omitted
complete transversal does not belong to $\mathcal D$. The code projection is
\begin{equation}\label{eq:code-projection}
 P_{\mathcal C}F=bg(z),\qquad g(z)=\E[bF\mid z].
\end{equation}
Put $w_s=a(n_s)$ and $Q(F)=\sum_s w_s|\widehat F(A_s)|$.
Restricting the full atom inequality and using signed parity averaging gives
\begin{equation}\label{eq:subcode-norm}
 Q(h_u)\le\|u\|_2,
 \qquad V\left(\sum_sc_s\chi_{A_s}\right)=\sum_sw_s|c_s|.
\end{equation}

\begin{lemma}[Stability after restricting the transversals]\label{lem:subcode-stability}
Let $\mathcal A$ be the retained balanced parameter rays. Then
\[
 \dist(u,\mathcal A)^2\le K_c\|u\|_2(\|u\|_2-Q(h_u)),
 \qquad K_c=2^{24}\ell^2R^5.
\]
The same bound holds for $\|h_u-P_{\mathcal D}h_u\|_2^2$.
\end{lemma}
\begin{proof}
Write $\varepsilon=\|u\|_2-Q(h_u)$ and $S_0=2^{16}R^3$.
A closest full-code ray $u_*$ satisfies
$\|u-u_*\|_2^2\le S_0\|u\|_2\varepsilon$.
If its transversal is retained, use $u_*$. Otherwise, for its omitted
transversal $A_t$ and weight $w_t=a(|A_t|)$,
\[
 w_t|\widehat h_u(A_t)|\le Q_{\rm full}(h_u)-Q(h_u)
 \le\varepsilon,
 \qquad w_t|\widehat h_{u_*}(A_t)|=\|u_*\|_2.
\]
ReLU is $1$-Lipschitz and $\E[\bar x\bar x^\top]=I$, so
\[
 \|u\|_2\le\varepsilon+(1+w_{\max})\|u-u_*\|_2,
 \qquad w_{\max}=\max_t a(|A_t|).
\]
Since $\varepsilon\le\|u\|_2$, squaring gives
\[
 \|u\|_2^2\le[2(1+w_{\max})^2S_0+2]\|u\|_2\varepsilon.
\]
Use the zero ray in this case. The elementary bounds $n\le a(n)\le3n$
give $w_{\max}\le3\ell(R+4)\le6\ell R$, so the bracket is at most
$K_c$. For the lower bound, use $a(2)\ge2$ and
$a(n+2)/a(n)=\sqrt{n(n+2)}/(n-1)\ge(n+2)/n$.
For the upper bound, induction gives
$\binom{2k}k/4^k\ge1/(2\sqrt k)$ for $k\ge1$, because
$(2k+1)^2\ge4k(k+1)$. Substituting $k=(n-2)/2$ gives
$a(n)\le2\sqrt{2n(n-2)}<3n$ for $n\ge4$; check $n=2$ directly.
Finally, the function generated by a retained ray belongs to
$\mathcal D$, and the same Lipschitz estimate proves function stability.
\end{proof}

For a representation of variation cost $B$ and any physical code dual
$\Lambda$ with $|\widehat\Lambda(A_s)|\le w_s$, write
$D=B-\langle\Lambda,F\rangle$. Weighted Cauchy--Schwarz gives
\begin{equation}\label{eq:network-stability}
 \|F-P_{\mathcal D}F\|_2^2\le K_cBD.
\end{equation}
Indeed the weighted sum of atom deficits is at most $D$. The statement
also holds for $B=V(F)$ by approximation.

\begin{lemma}[Observability outside independent syndrome sets]
\label{lem:observability}
If $U\subseteq G$ satisfies $U\cap(U+a_j)=\varnothing$ for every
column $a_j$, then every $A\in\mathcal D\cap\mathcal C^\perp$ obeys
\begin{equation}\label{eq:half-observability}
 \|A\|_2^2\le2\E[\one_{z\notin U}A^2].
\end{equation}
\end{lemma}
\begin{proof}
Each Fourier support in $A$ is a proper subset of a retained
transversal. Assign it to the first rail pair $j$ at which its
restriction is neither complete rail, and write $A=\sum_j A_j$.
Conditioned on $z$, two physical characters have nonzero covariance
only if their symmetric difference consists of full rail pairs with
incidence vector $(a_i\cdot s)_i$ for some $s\in G$.
At a defective pair, a proper subset of one rail cannot be complemented
within the full pair while remaining a subset of a single rail. Thus
the two restrictions must agree there. Comparing first defective pairs
shows that $A_j,A_k$ are conditionally orthogonal for $j\ne k$.
Every Fourier mode $s$ of $\ell_j(z)=\E[A_j^2\mid z]$ has
$a_j\cdot s=0$, so $\ell_j(z+a_j)=\ell_j(z)$.
The set $U$ contains at most one endpoint of each $a_j$-pair. Hence
$\E[\one_U\ell_j]\le\|A_j\|_2^2/2$. Summing proves the claim.
\end{proof}

\section{Exponentially many distinct completion costs}

Write $M=2^m$ and
$\widehat w(v)=M^{-1}\sum_s(-1)^{s\cdot v}w_s$.
Appendix~\ref{app:Bernstein} proves the exact representation
\begin{equation}\label{eq:Bernstein}
 a(r)=\sqrt{2\pi}(r-3/4)
 -2\sqrt{2\pi}\int_0^\infty e^{-rt/2}h(t)\,dt,
 \qquad 0<h(t)\le3.
\end{equation}

\begin{lemma}[Walsh signs and a high-weight bound]\label{lem:Walsh}
Every nonconstant $\widehat w(v)$ is strictly negative.
For $|v|\ge q\ge3$, with $k=\lceil q/2\rceil$,
\begin{equation}\label{eq:Walsh-tail}
 |\widehat w(v)|\le\frac{31}{\ell R}\left(\frac8R\right)^k.
\end{equation}
\end{lemma}
\begin{proof}
Write $n_s=n_0+\sum_j\delta_j(a_j\cdot s)$, where the Boolean
linear form is interpreted as zero or one and $\delta_j\in\{2,4\}$.
Each factor $e^{-t\delta_j(a_j\cdot s)}$ has Fourier coefficients
$(1+e^{-t\delta_j})/2$ at zero and
$(1-e^{-t\delta_j})/2$ at $a_j$. Therefore the Fourier coefficients
of $e^{-t(n_s-n_0)}$ form the probability law of a sum of independent
column-valued Bernoulli variables. Every syndrome has positive
probability for $t>0$, since the loop columns alone span $G$.
The affine part of \eqref{eq:Bernstein} has nonpositive nonconstant
coefficients, while its negative integral has strictly negative ones.
This proves the sign assertion.

Each column has Hamming weight at most two. A column subset summing to
$v$ therefore has at least $k$ members. Each Bernoulli success probability
is at most $\delta_jt/2\le2t$, so the union bound over $k$ successes
bounds the $v$-coefficient by $\binom\ell k(2t)^k$.
When $|v|\ge3$, the affine part has zero coefficient there.
Rescale the integral in \eqref{eq:Bernstein} by $t/2\mapsto t$.
Its density is bounded by $12\sqrt{2\pi}$, whence
\[
 |\widehat w(v)|\le12\sqrt{2\pi}\binom\ell k2^k
          \int_0^\infty e^{-n_0t}t^k\,dt
 \le\frac{31}{n_0}\left(\frac{2\ell}{n_0}\right)^k.
\]
This implies the stated, weaker bound.
\end{proof}

\begin{lemma}[Distinct scalar levels from a path transfer matrix]
\label{lem:path-levels}
The set $\{\widehat w(v):v\in G\}$ has at least $2^{m/2}$ distinct values.
\end{lemma}
\begin{proof}
Let $B$ be the vertex-edge incidence matrix of the path and put
$a=\tanh(2t)$, $e=\tanh t$. Expanding the Fourier factors just used gives
\[
 \widehat{e^{-tn_s}}(v)
 =e^{-tn_0}\left(\frac{1+e^{-4t}}2\right)^m
             \left(\frac{1+e^{-2t}}2\right)^{m-1} Z_v(a,e),
\]
\[
 Z_v(a,e)=\sum_{y\in\F_2^{m-1}}e^{|y|}a^{|v+By|}.
\]
Define
\[
 M_0=\begin{pmatrix}1&ae\\a&e\end{pmatrix},\qquad
 M_1=\begin{pmatrix}a&e\\1&ae\end{pmatrix},\qquad
 c_0=\binom1a,\qquad c_1=\binom a1.
\]
Successively summing the edge bits gives
\[
 \binom{Z_v}{Z_{v+e_1}}
 =M_{v_1}\cdots M_{v_{m-1}}c_{v_m}.
\]
Choose $t=\operatorname{arctanh}(1/2)$, so $e=1/2$ and $a=4/5$.
On the interval $I=[4/5,5/4]$, the ratio map of $M_0$ is
\[
 f(r)=\frac{r+ae}{ar+e},\qquad f(I)=[20/19,11/10]\subset(1,5/4].
\]
The ratio map of $M_1$ is $1/f(r)$ and lies below one in $I$.
Both are invertible because $\det M_0=e(1-a^2)>0$.
The terminal ratios are $1/a$ and $a$.
Thus $Z_v/Z_{v+e_1}$ determines $v$: the side of one identifies the
first bit, inversion removes that matrix, and iteration reaches the
distinct terminal ratios. There are $2^m$ distinct ordered pairs
$(Z_v,Z_{v+e_1})$. If there are $L$ scalar values among all $Z_v$,
there are at most $L^2$ ordered pairs, so $L\ge2^{m/2}$.

It remains to transfer this scalar separation to actual parity costs.
All even degrees belong to $[n_0,n_0+6m-2]$, whose width is less than
$2\sqrt{n_0}$. Two distinct integers in the same squarefree class,
$ru^2<rv^2$, both at least $n_0$, differ by at least
$r(2u+1)>2\sqrt{n_0}$. Hence all possible even degrees have distinct
squarefree parts. Each $a(n)$ is a nonzero rational multiple of
$\sqrt n$, so these costs are linearly independent over $\mathbb Q$.
For completeness, distinct squarefree radicals are independent because
the sign-changing automorphisms of the field generated by their prime
square roots separate the corresponding characters; averaging against
each character isolates its coefficient.

Consequently equality of two Walsh costs is equivalent to equality of
the signed counts at every degree, hence to equality of the polynomials
$\sum_s(-1)^{s\cdot v}x^{(n_s-n_0)/2}$.
Distinct transfer values at the chosen $t$ imply distinct polynomials
and therefore distinct Walsh costs. This proves the claim.
\end{proof}

\begin{lemma}[An observable affine set retaining many levels]
\label{lem:live-set}
For $q=q_m$ and all sufficiently large $m$, there is an affine coset
$U=v_0+H$ and a subset $\mathcal V\subseteq U$ of size $T_m$ such that
every $v\in U$ has $|v|\ge q$, the costs indexed by $\mathcal V$ are
pairwise distinct, and $U\cap(U+a_j)=\varnothing$ for every column.
The codimension of $H$ is between $q$ and $q+2$.
\end{lemma}
\begin{proof}
Take one representative of each distinct level in Lemma~\ref{lem:path-levels}.
Discard those of weight below $q$, at most $\sum_{j<q}\binom mj$ points.
Each remaining point contains a $q$-subset $J$ of its support. Pigeonhole
over $\binom mq$ such subsets and the four values of the two parities
with normal vectors
$\mathbf1=(1,\ldots,1)$ and $a_{\rm alt}=(1,0,1,0,\ldots)$.
Some affine set
\begin{equation}\label{eq:U}
 U=\{z:z_j=1\ (j\in J),\quad
       \mathbf1\cdot z=\alpha,\quad a_{\rm alt}\cdot z=\beta\}
\end{equation}
contains at least $T_m$ of these representatives. Select that many as
$\mathcal V$. Ties in all selections may be settled lexicographically.
The set is nonempty and excludes zero. Its defining normals have rank
between $q$ and $q+2$. A loop translation changes the all-ones parity;
an edge translation changes the alternating parity. Thus $U$ is
independent in every column direction.

Finally $q_m=O(m/\log m)$, so
$\log_2\binom m{q_m}=O(m\log\log m/\log m)=o(m)$ and
$\sum_{j<q_m}\binom mj=2^{o(m)}$.
The numerator defining $T_m$ is positive for all sufficiently large
$m$, and $T_m=2^{m/2-o(m)}$ as claimed.
\end{proof}

\section{Calibration that forces a single exceptional fiber}

Put $M=2^m$ and use $U=v_0+H$ and $\mathcal V$ from
Lemma~\ref{lem:live-set}. Write $K=|U|$, $r=\operatorname{codim}H$,
$J_0=H+\spanop(v_0)$, and $E=H^\perp\setminus J_0^\perp$.
Then $|E|=2^{r-1}\ge2$, and every character in $E$ equals $-1$ on $U$.

The target logical label is
\[
 t(z)=\begin{cases}+1,&z=0,\\-1,&z\ne0,\end{cases}
 \qquad f(x)=b(x)t(z(x)).
\]
Initially reveal every physical input whose logical coordinate lies
outside $\mathcal V$, including all fibers in $U\setminus\mathcal V$.
Subsequently whole fibers of $\mathcal V$ will be revealed in a specified order. Let $U_{\rm live}\subseteq\mathcal V$ be the nonempty set of
fibers not yet revealed. All known labels on a removed fiber are the
negative logical label. Write $\mathcal K=\{x:z(x)\notin U\}$ for the fixed calibration region. Extra revealed fibers impose additional constraints but are not included in averages with subscript $\mathrm{cal}$.

For $v\in U$, define the star score
\[
 g_v(z)=-\tfrac12+\one_{z=0}+\one_{z=v},\qquad
 F_v(x)=b(x)g_v(z(x)).
\]
It agrees with the target on every fiber except $v$, and has absolute
margin $1/2$ everywhere. Its code coefficients have negative sign at
$s=0$, are nonnegative elsewhere, and vanish on $E$.

Define the code dual through its logical Fourier coefficients:
\begin{equation}\label{eq:dual}
 \widehat\lambda_s=w_s-2w_0\one_{s=0}-p_E\one_{s\in E},
 \qquad p_E=\frac{2w_0}{|E|}.
\end{equation}
Its physical version is $b(x)\lambda(z(x))$. Since $w_s\ge w_0$,
$|\widehat\lambda_s|\le w_s$, so it is a unit variation-norm dual.
Fourier inversion gives
\begin{equation}\label{eq:dual-values}
 \lambda(z)=M\widehat w(z)-2w_0-2w_0\one_H(z)+2w_0\one_U(z).
\end{equation}
In particular $\lambda(v)=M\widehat w(v)$ on $U$, while on the fixed calibration region
its sign agrees with $t$ and its magnitude is at least $2w_0$.

The stars saturate the dual and have exact costs
\begin{equation}\label{eq:star-costs}
 V(F_v)=C_v=C_*+\widehat w(v),\qquad
 C_*=\frac1M\sum_sw_s+\frac{w_0}{2}-\frac{4w_0}{M}.
\end{equation}
For $v\in\mathcal V$, Lemma~\ref{lem:live-set} makes these costs
pairwise distinct. Only these representatives can remain live.

\subsection{A linear rounding estimate}

Consider any finite ReLU score $F$ with absolute margin at least $1/2$
on the full cube and correct margins on all known fibers. Let $B$ be its
variation cost, or its variation norm. We only need consider
$B\le\max_v C_v$, since the remaining stars are competitors. Write its
code projection as $bg$, and set
\[
 p_v=g(v)+\tfrac12\quad(v\in U),\qquad
 D=B-\langle b\lambda,F\rangle,
\]
\[
 S=\E_{\rm cal}\bigl[|\lambda(z)|(fF-\tfrac12)\bigr].
\]
The expectation with subscript ``cal'' is the original normalized
physical expectation with the indicator of $\mathcal K$. Both $D$ and
$S$ are nonnegative. A direct decomposition gives
\begin{equation}\label{eq:decomposition}
 B=C_*+D+S+\sum_{v\in U}\widehat w(v)p_v.
\end{equation}

\begin{lemma}[Physical variance control]\label{lem:variance}
For these scores,
\begin{equation}\label{eq:variance}
 L:=\|F-P_{\mathcal C}F\|_2^2
 \le C_{\rm phys}(D+S),\qquad C_{\rm phys}=2^{30}\ell^3R^6.
\end{equation}
\end{lemma}

\begin{proof}
On a fiber in the fixed calibration region put $e=fF-1/2\ge0$.
The signed conditional expectation \eqref{eq:code-projection} gives
\[
 \E_{\rm cal}|F-P_{\mathcal C}F|\le2\E_{\rm cal}e\le S/w_0.
\]
A variation-cost-$B$ representation satisfies
$\|F\|_\infty\le\sqrt{d+1}B$, and $P_{\mathcal C}$ is an
$L^\infty$ contraction. Hence the observed squared residual is at most
$4\sqrt{d+1}B\E_{\rm cal}e$.
Write $E_0=F-P_{\mathcal D}F$ and
$A=P_{\mathcal D}F-P_{\mathcal C}F$.
Lemma~\ref{lem:observability}, $(u-v)^2\le2u^2+2v^2$, and
\eqref{eq:network-stability} give
\[
 \|A\|_2^2\le16\sqrt{d+1}B\E_{\rm cal}e+4\|E_0\|_2^2,
 \qquad \|E_0\|_2^2\le K_cBD.
\]
The two residuals are orthogonal, so
\[
 L\le B\left[\frac{8\sqrt{d+1}}{w_0}S+5K_cD\right].
\]
We may use $B\le5\ell(R+4)\le10\ell R$, $w_0\ge\ell R$,
and $d+1\le4\ell R$. The coefficients are bounded by
$160\sqrt{\ell R}$ and $50\cdot2^{24}\ell^3R^6$, respectively.
Both are at most $C_{\rm phys}=2^{30}\ell^3R^6$.
\end{proof}

\begin{lemma}[Rounding to a live star]\label{lem:rounding}
For some $v\in U_{\rm live}$,
\begin{equation}\label{eq:rounding}
 \|p-e_v\|_1\le B_{\rm round}(D+S),\qquad
 B_{\rm round}=2^{33}2^m\ell^3R^6.
\end{equation}
Here $e_v$ is the point mass at $v$, indexed by all of $U$.
\end{lemma}

\begin{proof}
Put $d_z=g(z)-(-1/2+\one_{z=0})$ outside $U$. Calibration and
\eqref{eq:dual-values} imply
\[
 \sum_{z\notin U}|d_z|\le\frac{MS}{2w_0}.
\]
The Fourier norm inequality and dual slack give
\[
 \sum_{s\in E}|\widehat g(s)|\le D/p_E,
 \qquad 2w_s(-\widehat g(s))_+\le D\quad(s\notin E\cup\{0\}).
\]
Set $\eta=M(D+S)/(2w_0)$. For every $s\in E$,
\[
 M\widehat g(s)=1-\sum_{v\in U}p_v
                  +\sum_{z\notin U}d_z(-1)^{s\cdot z}.
\]
Average the absolute-value bound over $E$ and use
$|E|p_E=2w_0$. This gives
\[
 \left|\sum_{v\in U}p_v-1\right|\le\eta.
\]
For any nontrivial character of $H$, choose two extensions to $G$ having
opposite signs on $v_0$. They differ by an element of $E$ and both lie
outside $H^\perp$. Applying the negative-coefficient inequality to these
two extensions bounds the absolute value of the unnormalized Fourier
transform of $p$ by $1+\eta$. The zero frequency obeys the same bound.
Parseval on the affine cube $U$ gives
\begin{equation}\label{eq:p-l2}
 \sum_{v\in U}p_v^2\le(1+\eta)^2.
\end{equation}

Every value of the signed physical score $bF$ has absolute value at least
$1/2$. Its conditional mean on a fiber $v$ is $p_v-1/2$. Therefore
\[
 V_{\rm frac}:=\sum_{0<p_v<1}p_v(1-p_v)\le ML.
\]
Suppose first $\eta\le1$. Then
$\sum p_v^2-\sum p_v\le4\eta$, so
\[
 \sum_{p_v\notin[0,1]}(p_v^2-p_v)\le4\eta+V_{\rm frac}.
\]
Round every coordinate to its nearest $b_v\in\{0,1\}$. The distance is
at most $p_v^2-p_v$ outside $[0,1]$ and at most $2p_v(1-p_v)$ inside.
Consequently
\[
 \|p-b\|_1\le4\eta+3V_{\rm frac},\qquad
 \left|\sum b_v-1\right|\le5\eta+3V_{\rm frac}.
\]
A revealed whole fiber has $p_v\le0$, hence $b_v=0$. If $b$ is nonzero,
choose a live $v$ with $b_v=1$; if $b=0$, choose any live $v$. In both
cases $\|b-e_v\|_1=|\sum b_v-1|$, so
\[
 \|p-e_v\|_1\le9\eta+6V_{\rm frac}.
\]

For $\eta>1$, use the logical Fourier algebra norm
$\|h\|_{\rm A}=\sum_s|\widehat h(s)|$. Multiplication is submultiplicative
for this norm, and $\|\one_U\|_{\rm A}=1$. Extend $p$ by zero outside
$U$, obtaining $\widetilde p=(g+1/2)\one_U$. The atom norm bound gives
\[
 \|\widetilde p\|_{\rm A}
 \le\|g\|_{\rm A}+\tfrac12
 \le B/w_0+\tfrac12\le10.5.
\]
Let $P_h=\sum_{u\in H}p_{v_0+u}h(u)$ for characters $h$ of $H$.
Each such character has $M/K$ extensions to $G$, whose coefficients of
$\widetilde p$ have absolute value $|P_h|/M$. Consequently
\[
 \frac1K\sum_h|P_h|=\|\widetilde p\|_{\rm A}\le10.5,
 \qquad
 \sum_vp_v^2=\frac1K\sum_h|P_h|^2
 \le10.5(1+\eta)\le21\eta.
\]
Together with $\sum p_v\ge1-\eta$, this yields
$\sum(p_v^2-p_v)\le22\eta$. The same coordinate rounding now gives
\[
 \|p-b\|_1\le22\eta+3V_{\rm frac},\qquad
 \left|\sum b_v-1\right|\le23\eta+3V_{\rm frac}.
\]
Choosing a live point as above proves
$\|p-e_v\|_1\le45\eta+6V_{\rm frac}$. Thus in all cases
\[
 \dist_1(p,\{e_v:v\in U_{\rm live}\})
 \le M\left(\frac{45}{2w_0}+6C_{\rm phys}\right)(D+S).
\]
Using $w_0\ge\ell R$ and Lemma~\ref{lem:variance} bounds this constant by
$2^{33}M\ell^3R^6=B_{\rm round}$.
\end{proof}

\begin{proposition}[Every minimum-cost completion is the cheapest star]
\label{prop:classification}
At every stage with $U_{\rm live}\ne\varnothing$, the unique
minimum-variation-cost score consistent with all known whole fibers and
the closed full-cube margins is $F_v$, where $v$ is the cheapest live
fiber in the order \eqref{eq:star-costs}.
\end{proposition}

\begin{proof}
Let $W=\max_{v\in U}|\widehat w(v)|$. For a score no more expensive than
a live star, choose the live vertex furnished by
Lemma~\ref{lem:rounding}. Equations \eqref{eq:decomposition} and
\eqref{eq:rounding} imply
\[
 B\ge C_v+(1-WB_{\rm round})(D+S).
\]
The parameter choice makes $WB_{\rm round}<1/2$. Indeed
Lemma~\ref{lem:Walsh}, $R=2^{42}\ell$, and $q_m=2k_m$ give
\[
 \log_2(WB_{\rm round})
 <m+248+7\log_2\ell-k_m(39+\log_2\ell)\le-2.
\]
The last inequality is exactly the defining lower bound on $k_m$.

A cheapest star attains equality at its cost. Any other minimizer must
therefore have $D=S=0$. Lemma~\ref{lem:variance} then gives $L=0$, and
Lemma~\ref{lem:rounding} gives $p=e_v$. Calibration slack zero fixes all
remaining code values, so the physical score is exactly $F_v$. The star
costs are distinct, proving uniqueness.
\end{proof}

\section{Finite width, the input sequence, and the posterior}

Each star has an exact minimum-variation representation obtained by
summing signed parity projections. Its number of neurons is at most
\[
 \sum_{s\in G}2^{n_s}\le2^{m+n_{\max}}<N_0.
\]
At every $N\ge N_0$ it can therefore be zero-padded and balanced to attain
energy $2\sqrt d\,C_v$. Conversely
\eqref{eq:energy-variation} and Proposition~\ref{prop:classification}
show that every minimum-energy completion has precisely the cheapest
live star score.

Order $\mathcal V$ by increasing $C_v$. The input permutation first lists
all physical inputs outside $\mathcal V$. Then it lists the complete
physical fibers in this cost order, with any fixed order inside each fiber.
Designate the first input in each such fiber. Just before this input is
queried, the entire preceding fibers and all fibers outside $\mathcal V$ have
been revealed. The unique minimum-energy score is therefore the star that is
positive on this next fiber. The target is negative there. Multiplication
by the base character preserves this disagreement.

For completeness, the passage to the actual posterior uses only a
finite-dimensional small-noise argument. At a designated round, replace
all strict constraints by their closed versions. The feasible set is
closed and the quadratic energy is coercive, so minima exist. Requiring
the correct sign at the designated point excludes every energy minimizer.
It therefore raises the minimum energy by a strictly positive amount.
The strict feasible set has the same infimum as its closed version:
scale the readout of a minimizing star by $1+\varepsilon$ and let
$\varepsilon\downarrow0$.

Here is a direct Gaussian integral bound. Let $E_0$ be the minimum energy
and $E_1>E_0$ the minimum on the closed correct-prediction set; if the
latter set is empty, its posterior probability is zero. Choose
$a\in(E_0,E_1)$. A sufficiently slightly scaled star has a fixed open
parameter ball $\mathcal B$ that is strictly feasible and satisfies
$K_N\le a$ throughout the ball. Put $P=N(d+2)$ and choose $\rho\in(0,1)$
so that $(1-\rho)E_1>a$. The unnormalized correct-prediction integral is
at most
\[
 \int_{K_N\ge E_1}e^{-K_N/(2\sigma^2)}\,d\theta
 \le e^{-(1-\rho)E_1/(2\sigma^2)}
       \int_{\R^P}e^{-\rho K_N/(2\sigma^2)}\,d\theta
 =c_\rho\sigma^P e^{-(1-\rho)E_1/(2\sigma^2)},
\]
where $c_\rho$ is finite and independent of $\sigma$. The denominator is
at least $\operatorname{vol}(\mathcal B)e^{-a/(2\sigma^2)}$.
Their ratio tends to zero. The Gaussian normalizing constant and the
initial full-margin conditioning cancel from this ratio. This argument
sums all parameterizations, including non-minimizing ones; it does not
make a tie-breaking assumption.

There are finitely many designated rounds for fixed $m$. Choose the prior
scale small enough that their correct-label probabilities sum to less
than $\eta$. A union bound then proves the posterior-sampling claim, and
probabilities below $1/2$ and $2/3$ prove the other two claims.

It remains to verify the teacher resources. Compute each rail parity,
then the $m$ logical parities (each uses at most three rail-pair parities)
and their AND; multiplying by the base
parity gives $f$. This uses $O(d)$ bounded-fan-in Boolean gates. For a
shallow ReLU teacher use the score
\begin{equation}\label{eq:teacher}
 F_{\rm teacher}(x)=\sum_{i=1}^m\chi_{A_{e_i}}(x)
                   -(m-1)\chi_{A_0}(x)
 =b(x)\left(\sum_i(-1)^{z_i}-(m-1)\right).
\end{equation}
It has margin one and the target sign. Parity on $r$ bits is a
piecewise-linear function of their sum on the $r+1$ attainable values.
Linear interpolation represents it by at most $r+1$ affine ReLU neurons,
with incoming coordinate weights one, biases bounded by $r$, and readout
coefficients bounded by two. The constant term is a neuron on the
constant input. Applying this to the $m+1$ parities in
\eqref{eq:teacher} gives width at most
\[
 (n_0+1)+\sum_{i=1}^m(n_0+5+2\deg_{\rm path}(i))
 =(m+1)(n_0+1)+8m-4=C_m.
\]
All incoming weights have magnitude at most one, biases have magnitude
at most $n_0+8$, and readout coefficients have magnitude at
most $2(m-1)$. Thus every coefficient is bounded by $C_m$. This completes
the proof of Theorem~\ref{thm:main}.

\section{The same lower bound for standard softplus}
\label{sec:softplus}

The conclusion also holds for a smooth activation. Let
\[
 \phi(t)=\log(1+e^t).
\]
Replace ReLU by $\phi$ in \eqref{eq:network}, leaving the Gaussian prior,
pre-data full-cube margin conditioning, and exact-label likelihood
unchanged. In particular the quadratic energy remains \eqref{eq:energy}.

\begin{theorem}[Softplus at every sufficiently large width]
\label{thm:softplus}
Use the target, input permutation, and designated rounds in
Theorem~\ref{thm:main}. For each sufficiently large $m$ there is a finite threshold
$N_\phi(m)$ such that, for every fixed $N\ge N_\phi(m)$, the standard
softplus learner's correct-label posterior probability tends to zero at
every designated round as $\sigma\downarrow0$.
The target has a one-hidden-layer softplus implementation of width at
most $C_m$, output margin greater than one, and every coefficient bounded
by $24C_m$. Equivalently, padding to width $\widetilde C_m=24C_m$ gives
both width and coefficient bounds $\widetilde C_m=\Theta(m^3)$.
Thus the cumulative mistake lower bound remains
$\exp(\Omega(\widetilde C_m^{1/3}))$.
The three prediction consequences in Theorem~\ref{thm:main} hold for
this learner as well.
\end{theorem}

The threshold in this extension is existential; the explicit ReLU
threshold $N_0$ need not suffice. The theorem concerns standard softplus,
not the bounded activation tanh.

\subsection{A bias-independent atom estimate}

For a physical weight vector $w\in\R^d$, excluding the constant-input
coefficient, define
\[
 G(w)=\frac1\pi\sum_s a(n_s)\int_0^\infty
 \prod_{j\in A_s}|\sin(w_jt)|
 \prod_{j\notin A_s}|\cos(w_jt)|\,\frac{dt}{t^2}.
\]
The coefficient integral \eqref{eq:Fourier-integral} implies
\[
 Q\bigl(\relu(b+w\cdot x)\bigr)\le G(w)\le\|w\|_2
 \qquad\text{for every }b\in\R.
 \tag{S1}
\]
To justify the second inequality, first sum over all $\ell$-pair
transversals to obtain $G_{\rm full}$. On nonnegative block-constant
vectors $V$, this equals the full zero-bias coefficient norm.
The generalized H\"older sampling argument in Appendix~\ref{app:rail}
therefore gives $G_{\rm full}(w)\le\|w\|_2$; retaining only selected
nonnegative terms gives $G(w)\le G_{\rm full}(w)$.

Let $\mathcal A$ be the retained balanced rays in physical weight
space. The full-code balanced stability and sampling argument in
Appendix~\ref{app:rail} applies to $G_{\rm full}$ without a bias
coordinate and gives squared-distance constant $2^{16}R^3$.
The restriction argument in Lemma~\ref{lem:subcode-stability} then
applies to $G$ as well. Indeed, if the closest full ray has omitted
transversal $A_t$, its omitted integral satisfies
\[
 w_t|\widehat{\relu(w\cdot x)}(A_t)|
 \le G_t(w)\le G_{\rm full}(w)-G(w)\le\|w\|_2-G(w).
\]
The same Lipschitz coefficient comparison permits using the zero ray
in that case. Consequently
\[
 \dist(w,\mathcal A)^2\le K_c\|w\|_2(\|w\|_2-G(w)).
 \tag{S2}
\]
This statement places no restriction on the bias. If $w_*\in\mathcal A$,
then any function of $b+w_*\cdot x$ belongs to the downward Fourier
space $\mathcal D$, because it depends only on one selected transversal.

For a centered logistic random variable $Z$,
\[
 \phi(t)=\E\relu(t+Z).
 \tag{S3}
\]
Both sides have derivative $(1+e^{-t})^{-1}$ and tend to zero as
$t\to-\infty$. Averaging (S1) proves
$Q(\phi(b+w\cdot x))\le G(w)\le\|w\|_2$.
Softplus is $1$-Lipschitz, so comparison with a closest $w_*\in\mathcal A$
while preserving the bias, together with $\E[xx^\top]=I$, gives
\begin{align}
 \|\phi(b+w\cdot x)-P_{\mathcal D}\phi(b+w\cdot x)\|_2^2
 &\le K_c\|w\|_2\bigl(\|w\|_2-G(w)\bigr)\notag\\
 &\le K_c\|w\|_2
       \bigl(\|w\|_2-Q(\phi(b+w\cdot x))\bigr).
 \tag{S4}
\end{align}
The constant component created by $\phi(0)=\log2$ lies in $\mathcal D$.

\subsection{A uniform energy gap}

For a softplus network put $B=K_N/(2\sqrt d)$. The arithmetic--geometric
mean inequality gives
\[
 B\ge\sum_j|v_j|\sqrt{b_j^2+\|w_j\|_2^2}
 \ge\sum_j|v_j|\|w_j\|_2.
\]
For the stage's code dual $b(x)\lambda(z(x))$, write
$D=B-\langle b\lambda,F\rangle$. The atom bound and weighted
Cauchy--Schwarz applied to (S4) show
\[
 Q(F)\le B,\qquad D\ge0,\qquad
 \|F-P_{\mathcal D}F\|_2^2\le K_cBD.
 \tag{S5}
\]
Since $|\phi(t)|\le|t|+\log2$ and
$\sum_j|v_j|\le\sqrt{N\sum_jv_j^2}\le\sqrt{K_N}$, we also have the
width-independent estimate
\[
 \|F\|_\infty
 \le\sqrt{d+1}\,B+(\log2)\sqrt{2\sqrt d\,B}.
 \tag{S6}
\]

Set $B_{\rm cap}=5\ell(R+4)\le10\ell R$. Every star cost is at most
$4.5\ell(R+4)$, so this cap exceeds every star cost by more than one.
For $B\le B_{\rm cap}$, (S6) is bounded by
\[
 H_{\rm cap}=3\sqrt{d+1}\,B_{\rm cap}.
\]
Repeat the proof of Lemma~\ref{lem:variance} using (S5) and this uniform
supremum bound. It gives
\[
 L\le\frac{8H_{\rm cap}}{w_0}S
       +5\cdot K_cB_{\rm cap}D.
\]
The two coefficients are at most $480\sqrt{\ell R}$ and
$50\cdot2^{24}\ell^3R^6$, respectively, so the same
$C_{\rm phys}=2^{30}\ell^3R^6$ applies. The calibrated decomposition,
Fourier slack inequalities, and rounding lemma therefore hold with
exactly the same $B_{\rm round}$ and absorption estimate.

To make the uniformity in width explicit, consider pairs $(B,F)$ where
$F$ is an arbitrary vector of real outputs indexed by the finite input
cube. Impose $0\le B\le B_{\rm cap}$,
$\|F\|_\infty\le H_{\rm cap}$, the closed full-cube and history margins,
and the three inequalities (S5). This is a compact set in a
finite-dimensional space. Every softplus network satisfying the closed
full-cube and history margins below the energy cap belongs to it, and
every ReLU star paired with its exact variation cost
also belongs to it.

Only these inequalities are used by the calibrated decomposition,
variance control, and rounding argument. Consequently every pair in
this compact set satisfies
\[
 B\ge C_v+(1-WB_{\rm round})(D+S)
\]
for some live star $v$. In particular $B\ge C_{\min}$, with equality
only for the unique cheapest live star $F_{\min}$.
Requiring the correct closed margin at the designated next point excludes
that score. Compactness then supplies a strictly positive gap $g$ above
$C_{\min}$ on the correct-prediction subset, if that subset is nonempty.
Energies above $B_{\rm cap}$ are separated by more than one already.
Thus, for some $g'>0$ independent of learner width, every softplus network
satisfying those margins and predicting correctly at this point has
\[
 K_N\ge2\sqrt d\,(C_{\min}+g').
 \tag{S7}
\]

\subsection{Approaching the star energy at large width}

Take a finite ReLU representation of $F_{\min}$ attaining its exact
variation cost. Balance its parameters $(a_j,c_j)$ so that
\[
 d\sum_j\|a_j\|_2^2+\sum_jc_j^2=2\sqrt d\,C_{\min}.
\]
For every sufficiently large $N$, use softplus incoming vectors
$\sqrt N a_j$, readouts $c_j/\sqrt N$, and zero readouts on the remaining
neurons. Its output is
\[
 F_N(x)=\sum_j c_j\frac{\phi(\sqrt N a_j\cdot\bar x)}{\sqrt N}.
\]
Because $0\le\phi(t)-\relu(t)\le\log2$, these outputs converge uniformly
to $F_{\min}$, and their energy before a margin correction is exactly
$2\sqrt d\,C_{\min}$. Let $\delta_N>0$ bound the uniform error, with
$\delta_N\to0$. Multiplying the readouts by $1+4\delta_N$ gives strict
margins for large $N$, since
\[
 (1+4\delta_N)(1/2-\delta_N)>1/2
 \qquad(0<\delta_N<1/4).
\]
The energy still converges to $2\sqrt d\,C_{\min}$. It is therefore below
the correct-prediction lower bound (S7) for every sufficiently large
width. An open parameter ball around this strictly feasible witness and
the Gaussian integral comparison from the ReLU proof show that the
correct-label posterior probability tends to zero as $\sigma\downarrow0$
at each such fixed width. There are finitely many designated stages, so
one common threshold $N_\phi(m)$ works for all of them.

\subsection{A teacher with linearly bounded coefficients}

For even $r$, the parity spline in the ReLU teacher has hinge locations
on the grid $-r,-r+2,\ldots,r$, hinge coefficients of magnitude at most
two, and one constant neuron. At an attainable input sum, at most one
hinge argument is zero. The remaining nonzero hinge arguments have
absolute values $2,4,\ldots$, with at most two at each distance.
Replacing each ReLU term by $\phi(Gt)/G$ therefore changes the parity
value by at most
\[
 \frac{2\log2+4/(e^{2G}-1)+e^{-G}}G<\frac3G
 \qquad(G\ge1).
\]
Here the constant neuron contributes at most $e^{-G}/G$.
The logical coefficients in \eqref{eq:teacher} have total absolute value
$2m-1$, so choosing $G=24m$ makes the target-score error less than $1/4$.
Multiply the resulting readout by two. Its margin is then greater than
one on every Boolean input.

The hidden width remains $C_m$. All incoming coordinate weights have
magnitude at most $24m$, and biases have magnitude at most
$24m(n_0+8)\le24C_m$, since
$C_m-m(n_0+8)=n_0+m-3>0$.
The final readout coefficients have magnitude at most $4m/G<1$.
Padding to width $24C_m$ with zero readouts preserves the function and
puts both width and coefficient bounds under the same parameter. This
completes the proof of Theorem~\ref{thm:softplus}.

\appendix
\section{Proof of the exact unequal-rail norm}\label{app:rail}

Throughout this appendix, $m$ denotes the number of rail pairs in the
full-code theorem (called $p$ in Theorem~\ref{thm:rail}, and equal to
$\ell$ in the main construction). We use $A_s$,
$Q$, and $\mathcal D$ for the full transversals, full coefficient norm,
and full downward space. Thus the hypothesis is $R\ge2^{40}m\Delta$.
We also write $n_s=|A_s|$, $w_s=a(n_s)$,
$n_{\max}=\max_s n_s$, and $a_{\max}=\max_s a(n_s)$ locally.
These conventions do not change the subcode definitions in the main text.

\subsection{Two elementary analytic tools}

\begin{lemma}[Minimum-coordinate parity bound]\label{lem:mincoord}
For even $n\ge2$ and any affine ReLU ridge on $n$ Boolean coordinates,
\[
 |\widehat h([n])|\le\kappa_n\min_i|w_i|.
\]
\end{lemma}
\begin{proof}
A coordinatewise nondecreasing $a:\{-1,1\}^k\to[0,1]$ has top Fourier
coefficient at most
$2^{-k}\binom{k-1}{\lfloor(k-1)/2\rfloor}$ in absolute value.
To see this, average over coordinate permutations and express the resulting
nondecreasing sequence on Hamming levels as a convex combination of level
thresholds and a constant. The alternating binomial sum of a threshold
starting at level $j$ is
$(-1)^{k-j}\binom{k-1}{j-1}$.

Input flips make all incoming weights nonnegative. Integrate out a chosen
coordinate $i$. The difference
$[\relu(t+w_i)-\relu(t-w_i)]/(2w_i)$ lies in $[0,1]$ and is nondecreasing
in the other coordinates. Apply the preceding bound with $k=n-1$.
A zero weight gives a zero coefficient directly.
\end{proof}

The affine part of ReLU has no even Fourier coefficient of degree at least
two. The identity
\[
 |z|=\frac2\pi\int_0^\infty\frac{1-\cos(tz)}{t^2}\,dt
\]
therefore gives, up to its fixed degree sign, the coefficient integral
\begin{equation}\label{eq:Fourier-integral}
 \widehat h(A)=\frac{(-1)^{|A|/2+1}}\pi
 \int_0^\infty\cos(bt)
 \prod_{j\in A}\sin(w_jt)\prod_{j\notin A}\cos(w_jt)\,\frac{dt}{t^2}.
\end{equation}
For the even block supports under consideration these integrals are
absolutely convergent.

\begin{lemma}[Block balancing]\label{lem:balancing}
For any nonnegative weighted sum of absolute coefficients on nonempty
unions of disjoint even blocks, the supremum of its ratio to the augmented
incoming Euclidean norm is unchanged by restricting to zero bias and
nonnegative weights constant within each block.
\end{lemma}
\begin{proof}
Take absolute values in \eqref{eq:Fourier-integral}, discard bias and
irrelevant-coordinate cosine factors, and flip signs of incoming weights.
Call the resulting nonnegative integral $G(w)/\pi$. For block-constant
weights it equals the coefficient sum exactly, because every sine and
cosine exponent is even.

For a block of size $r$, generalized H\"older, on the disjoint union of
coefficient-index copies of the integration measure, gives
\[
 G(w)\le\prod_{j\text{ in block}}
 G(w\text{ with this whole block replaced by }w_j)^{1/r}.
\]
Iterate over blocks. If $V$ replaces each block by an independent uniformly
sampled coordinate magnitude from that block, the result is
$G(w)/\pi\le\exp(\E\log Q(h_V))$. If $M_*$ is the restricted supremum,
this is at most
$M_*\exp(\E\log\|V\|)\le M_*\sqrt{\E\|V\|^2}=M_*\|w\|$.
Bias and irrelevant weights only increase the original norm.
\end{proof}

\subsection{A quantitative H\"older estimate}

\begin{lemma}[Separation of parity densities]\label{lem:holder-gap}
Let $m\ge2$, let $R$ be positive and even, and put $n=mR$. For $c>0$,
define the probability density on $t>0$
\[
 p_c(t)=\frac{\sin(ct)^n}{\pi\kappa_n c t^2}=c p_1(ct).
\]
If $c_i\in[h/2,2h]$, set
\[
 \mathcal F(c)=\int_0^\infty\prod_{i=1}^m p_{c_i}(t)^{1/m}\,dt,
 \qquad
 Z(c)=\frac R{h^2}\sum_i(c_i-\bar c)^2,
 \quad \bar c=\frac1m\sum_i c_i.
\]
Then
\begin{equation}\label{eq:holder-gap}
 1-\mathcal F(c)\ge2^{-13}\min\{Z(c),1/m\}.
\end{equation}
Moreover, the weighted top-parity coefficient of an equal-block ridge is
\begin{equation}\label{eq:holder-parity}
 \mathcal Q(c):=a(n)\left|
 \widehat{\relu\!\left(\sum_i c_i\sum_{j=1}^R X_{ij}\right)}([n])
 \right|
 =\sqrt n\Bigl(\prod_i c_i\Bigr)^{1/m}\mathcal F(c)
 \le\sqrt{R\sum_i c_i^2}\,\mathcal F(c).
\end{equation}
\end{lemma}
\begin{proof}
The equal-weight coefficient integral \eqref{eq:Fourier-integral} gives
$\int_0^\infty\sin(ct)^n t^{-2}\,dt=\pi\kappa_n c$, so $p_c$ is
normalized. We first prove, for $c,d\in[h/2,2h]$, the squared Hellinger
separation
\begin{equation}\label{eq:hellinger-parity}
 \int_0^\infty(\sqrt{p_c}-\sqrt{p_d})^2\,dt
 \ge2^{-10}\min\{n(c-d)^2/h^2,1\}.
\end{equation}
Assume $d\ge c$ and compare their cumulative probabilities at
$t_0=\pi/(2\sqrt{cd})$. Their difference is the integral of $p_1$
between $(\pi/2)\sqrt{c/d}$ and $(\pi/2)\sqrt{d/c}$. This interval
contains $[\pi/2-L,\pi/2+L]$, where
\[
 L=\min\left\{\frac\pi8\frac{d-c}{h},\frac1{2\sqrt n}\right\}.
\]
Indeed its smaller distance from $\pi/2$ is
$\pi(d-c)/(2\sqrt d(\sqrt d+\sqrt c))\ge\pi(d-c)/(8h)$.
On the inner interval, Bernoulli's inequality gives $\sin(t)^n\ge7/8$,
while $t^2\le4$ and $\kappa_n\le1/\sqrt n$. Thus
$p_1(t)\ge\sqrt n/16$. The cumulative-probability difference is at
least
\[
 \frac{\sqrt n L}{8}\ge\frac1{32}
       \min\{\sqrt n|c-d|/h,1\}.
\]
Total variation dominates this difference, and its square is bounded by
the squared Hellinger distance by Cauchy--Schwarz. This proves
\eqref{eq:hellinger-parity}.

Apply arithmetic--geometric mean to the square roots of the densities.
Pointwise, their mean square minus the square of their mean is
$\frac1{2m^2}\sum_{i,j}(\sqrt{p_{c_i}}-\sqrt{p_{c_j}})^2$.
After integration this gives
\[
 1-\mathcal F(c)\ge\frac1{2m^2}
 \sum_{i,j}\int(\sqrt{p_{c_i}}-\sqrt{p_{c_j}})^2.
\]
For arbitrary real $x_i$, the elementary truncation estimate is
\[
 \frac1{m^2}\sum_{i,j}\min\{(x_i-x_j)^2,1\}
 \ge\frac14\min\left\{\frac1m\sum_i(x_i-\bar x)^2,\frac1m\right\}.
\]
If their range is at most one, the left side is twice their variance.
Otherwise one extreme point is at distance at least $1/2$ from at least
half of the points; counting both orders of these pairs gives a lower
bound $1/(4m)$. Apply this estimate to $x_i=\sqrt n c_i/h$ and use
\eqref{eq:hellinger-parity} to obtain \eqref{eq:holder-gap}.
Finally the coefficient integral gives the identity in
\eqref{eq:holder-parity}; arithmetic--geometric mean applied to $c_i^2$
gives its inequality.
\end{proof}

\subsection{Localization and the H\"older bootstrap}

By Lemma~\ref{lem:balancing}, first consider balanced nonnegative rails.
Within each pair call their magnitudes $a_i\ge b_i\ge0$, exchanging rail
names if necessary. Normalize the physical norm to one. Put $n=mR$ and
\[
 \alpha=\frac{a(m(R+\Delta))}{a(mR)}\le1+2\Delta/R.
\]
The bound follows from
$a(n+2)/a(n)=\sqrt{n(n+2)}/(n-1)$, which makes
$a(n)/(n-1)$ decreasing along even integers. Reducing every even
sine/cosine exponent to $R$ increases the nonnegative integrands. Hence
$Q_r\le\alpha Q_R$.

Clarkson's elementary inequality gives, for even $R$,
\[
 \sin^R(at)\cos^R(bt)+\cos^R(at)\sin^R(bt)
 \le\tfrac12\{\sin^R((a+b)t)+\sin^R((a-b)t)\}.
\]
Multiply over rail pairs, integrate, and apply
Lemma~\ref{lem:mincoord}. With independent symmetric signs $\varepsilon_i$,
\begin{equation}\label{eq:min-localization}
 Q_R\le\sqrt n\,\E\min_i(a_i+\varepsilon_i b_i).
\end{equation}
Suppose $Q_r\ge1-\varepsilon$, where
$0\le\varepsilon\le\Delta/(100R)$. Put
\[
 \beta=\frac\alpha{1-\varepsilon}\le1+3\Delta/R,
 \qquad h=(\beta\sqrt n)^{-1},\qquad y_i=a_i/h-1.
\]
The expectation in \eqref{eq:min-localization} is at least $h$ and at
most every $a_i$, so $y_i\ge0$. The physical norm gives
\[
 \sum_i(2y_i+y_i^2)\le m(\beta^2-1),\qquad
 s_0:=\sum_i y_i\le4m\Delta/R.
\]
Here and below \eqref{eq:large-R} implies $\Delta/R\le1/20$.
For $i\ne j$, concavity of the minimum gives
\[
 \E\min_k(a_k+\varepsilon_kb_k)
 \le\frac{a_i+a_j-\max(|a_i-a_j|,b_i)}2,
\]
so $b_i/h\le y_i+y_j\le s_0$ for every $j\ne i$. Write
$\zeta=\Delta/R\le2^{-40}/m$ and let
$c_i=a_i+\varepsilon_i b_i$. Since $s_0\le4m\zeta$, every such $c_i$
belongs to $[h/2,2h]$.

The same Clarkson reduction, before applying the minimum-coordinate
bound, gives $Q_r\le\alpha\E\mathcal Q(c)$, with $\mathcal Q$ from
Lemma~\ref{lem:holder-gap}. Put $\mathcal N(c)=\sqrt{R\sum_i c_i^2}$.
We have
\[
 \E\mathcal N(c)\le\sqrt{\E\mathcal N(c)^2}
 =\sqrt{R\sum_i(a_i^2+b_i^2)}\le1,
 \qquad \mathcal N(c)\ge\frac{\sqrt n h}{2}
 =\frac1{2\beta}\ge\frac14.
\]
By \eqref{eq:holder-parity} and $Q_r\ge1-\varepsilon$,
\[
 \E\bigl[\mathcal N(c)(1-\mathcal F(c))\bigr]
 \le1-\frac{1-\varepsilon}{\alpha}
 =1-\frac1\beta\le3\zeta.
\]
It follows that $\E(1-\mathcal F(c))\le12\zeta$, so
\eqref{eq:holder-gap} implies
\begin{equation}\label{eq:truncated-spread}
 \E\min\{Z(c),1/m\}\le2^{17}\zeta.
\end{equation}

Let $P$ denote orthogonal projection off the constant vector in
$\R^m$, and write $\mathbf a=(a_i)$ and $\mathbf b=(b_i)$. Then
\[
 \mu:=\E Z(c)=\frac R{h^2}
 \left(\|P\mathbf a\|^2+(1-1/m)\|\mathbf b\|^2\right),
 \qquad \E Z(c)^2\le3\mu^2.
\]
For the second assertion, a deterministic vector plus a Hilbert-space
Rademacher sum has fourth moment at most $A^2+6AB+3B^2\le3(A+B)^2$,
where $A$ is the deterministic squared norm and $B$ the expected squared
norm of the random sum. This follows by expansion: odd terms vanish,
and the centered fourth moment is at most $3B^2$.
If $\mu>0$, Cauchy--Schwarz gives
\[
 \Prb(Z(c)\ge\mu/2)\ge\frac1{12},\qquad
 \E\min\{Z(c),1/m\}\ge\frac1{24}\min\{\mu,1/m\}.
\]
Indeed $\E[Z\one_{Z\ge\mu/2}]\ge\mu/2$ and
$(\E[Z\one_{Z\ge\mu/2}])^2\le\E Z^2\Prb(Z\ge\mu/2)$.
The same resulting bound is immediate if $\mu=0$. Combining with
\eqref{eq:truncated-spread} yields
$\min\{\mu,1/m\}\le2^{22}\zeta<1/m$, hence $\mu\le2^{22}\zeta$.

Let $A$ be the selected transversal, $D_0=|A|$, and let $a$ be the
weighted mean of the selected rail magnitudes. Write the selected weights
as $a+\delta_j$, with $\sum_A\delta_j=0$, and denote the unused weights
by $u$. Set $\tau^2=\|\delta\|^2+\|u\|^2$. The weighted mean minimizes
the selected squared deviations. Since every rail size is at most $2R$,
\[
 \tau^2\le2R(\|P\mathbf a\|^2+\|\mathbf b\|^2)
 \le4h^2\mu.
\]
Since $a\ge h$, this proves the required bootstrap:
\begin{equation}\label{eq:tau-localization}
 \frac{\tau^2}{a^2}\le4\mu\le2^{24}\frac\Delta R
 \le\frac{2^{-16}}m<\frac1{256}.
\end{equation}

\subsection{Local concentration under parity conditioning}

\begin{lemma}[A local bound for tilted block sums]\label{lem:local-parity}
Divide independent signs into $m\ge2$ blocks of sizes $r_i\ge R$, where
$R\ge256$ and $R\ge10\log(m+1)$. Within block $i$, let the signs have
common mean $q_i$, with $|q_i|\le1/2$, and condition their product to an
arbitrary prescribed sign. If $S$ is the total sum and $D=\sum_i r_i$,
then
\[
 \sup_s\Prb(S=s)\le\frac{20}{\sqrt D}.
\]
\end{lemma}
\begin{proof}
A block of size $r$, mean $q$, and prescribed parity $\epsilon$ has
characteristic function
\[
 \frac{(\cos t+iq\sin t)^r+
       \epsilon(q\cos t+i\sin t)^r}{1+\epsilon q^r}.
\]
For $|t|\le\pi/8$, the ratio of the second numerator base's modulus to
the first is at most $2/3$, since its square is at most
\[
 \frac{1/4+(3/4)\sin^2(\pi/8)}{\cos^2(\pi/8)}<\frac49.
\]
Also $|\cos t+iq\sin t|\le e^{-t^2/8}$, by
$1-q^2\ge3/4$ and $\sin|t|\ge2|t|/\pi$. The block characteristic
modulus is therefore at most
\[
 \exp\bigl(-rt^2/8+4(2/3)^r\bigr).
\]
The product over blocks is at most $2e^{-Dt^2/8}$, because
$4m(2/3)^R<\log2$ under the hypotheses.

For $\pi/8\le|t|\le\pi/4$, both numerator bases have modulus at most
$e^{-1/20}$. Each block characteristic modulus is then at most
$3e^{-r/20}$; their product is at most $e^{-D/40}$ because $D\ge256m$.
Each conditioned block sum has a fixed residue modulo four, so the
modulus of the total characteristic function $\phi$ has period $\pi/2$.
Fourier inversion gives
\[
 \Prb(S=s)\le\frac2\pi\int_{-\pi/4}^{\pi/4}|\phi(t)|\,dt
 \le\frac{4\sqrt{8/\pi}}{\sqrt D}+e^{-D/40}
 <\frac{20}{\sqrt D}.
\]
\end{proof}

\begin{lemma}[The crossing remainder on each logical fiber]
\label{lem:local-crossing}
Consider uniform signs in disjoint even blocks of sizes at least $R$,
under arbitrary conditioning of their separate block parities. Select
$m\ge2$ of the blocks, with total size $D$, and let $S_0$ be their coordinate
sum. Let $T$ be a linear form with coefficients constant inside every
block, whose squared coefficient norm is $\tau^2>0$. Suppose
$R\ge256$, $R\ge10\log(m+1)$, and $a/\tau\ge16$. Define
\[
 \mathcal R=\relu(aS_0+T)-\relu(aS_0)-\one_{S_0>0}T
                         -\one_{S_0=0}\relu(T).
\]
Then, uniformly in the prescribed block parities,
\begin{equation}\label{eq:remainder}
 \E|\mathcal R|\le\frac{2^{14}e^{-192}\tau}{\sqrt D}.
\end{equation}
The same bound holds after conditioning on any logical fiber.
\end{lemma}
\begin{proof}
For one even block with coefficient $c$, parity conditioning gives
\[
 \E e^{tc\sum_j X_j}
 =\cosh(tc)^r[1+\epsilon\tanh(tc)^r]
 \le e^{rt^2c^2/2}(1+\tanh(tc)^2)
 \le e^{rt^2c^2}.
\]
Independence between the conditioned blocks therefore gives
$\E e^{tT}\le e^{t^2\tau^2}$ for every real $t$.

Put $\rho=a/\tau\ge16$ and take $t=8/\tau$. Exponential tilting by
$e^{\pm tT}$ preserves independence between parity-conditioned blocks.
In each selected block its underlying sign mean becomes
$\tanh(\pm tc)$, of absolute value at most
$t|c|\le8/\sqrt R\le1/2$. Apply Lemma~\ref{lem:local-parity} under the
tilted law and the moment bound to its normalizing constant. This yields
\[
 \E[e^{\pm tT}\one_{S_0=s}]\le\frac{20e^{64}}{\sqrt D}.
\]
Pointwise,
\[
 |\mathcal R|\le|T|\one_{S_0\ne0,\,|T|\ge a|S_0|}.
\]
The selected block sizes are even, so $S_0$ is even. Sum over its values
$\pm2\ell$, $\ell\ge1$, and use
$z\one_{z\ge u}\le(u+1/t)e^{t(z-u)}$ for $z,u\ge0$, together with
$e^{t|T|}\le e^{tT}+e^{-tT}$. We obtain
\begin{align*}
 \E|\mathcal R|
 &\le\frac{80e^{64}}{\sqrt D}
       \sum_{\ell\ge1}(2a\ell+\tau/8)e^{-16\rho\ell}\\
 &\le\frac{80\tau}{\sqrt D}(8\rho+1/4)e^{64-16\rho}
 \le\frac{2^{14}e^{-192}\tau}{\sqrt D}.
\end{align*}
For the second inequality use the geometric series with
$e^{-16\rho}\le1/2$; for the last, the preceding expression decreases
for $\rho\ge16$. Conditioning on a logical fiber is a mixture of the
separate block-parity conditionings, so the uniform bound survives.
\end{proof}

\subsection{A uniform cusp near a selected parity}

We first record a useful gap. Let $U=\sum_i c_i\sum_{j\in B_i}Y_j$ be a
sum of independent even blocks, each of size at least four and with a
constant coefficient inside a block. For every offset $b$,
\begin{equation}\label{eq:outside-gap}
 \sum_{\varnothing\ne I}\left|
 \widehat{|b+U|}\Bigl(\bigcup_{i\in I}B_i\Bigr)\right|
 \le\tfrac12\E|U|.
\end{equation}
Indeed \eqref{eq:Fourier-integral} bounds the coefficients by those at
$b=0$. Condition each block parity to be $(-1)^{|B_i|/2}$. The projected
value on this fiber equals the mean of $|U|$ minus the sum of the absolute
nonconstant block coefficients. This conditional distribution keeps the
mean zero and the variance unchanged. Its fourth moment is at most four
times the square of its variance: for a size-four block this follows by
direct expansion, and for larger blocks all fourth moments are unchanged.
Independence preserves the bound. H\"older therefore gives
$\E_{\rm fiber}|U|\ge\frac12\sqrt{\E U^2}\ge\frac12\E|U|$, proving
\eqref{eq:outside-gap}.

Write $S_0=\sum_{j\in A}X_j$ and $T=\delta\cdot X+u\cdot Y$. Then
\[
 \relu(aS_0+T)=\relu(aS_0)+\one_{S_0>0}T
              +\one_{S_0=0}\relu(T)+\mathcal R,
\]
where $\mathcal R$ is the crossing remainder from
Lemma~\ref{lem:local-crossing}. Its assumptions follow from
\eqref{eq:large-R} and \eqref{eq:tau-localization} whenever $\tau>0$.
Let $p_0=\Prb(S_0=0)$ and $\epsilon_0=(-1)^{D_0/2+1}$. The selected
coefficient satisfies
\[
 \epsilon_0\widehat h(A)=\kappa_{D_0}a
 -\frac{p_0}2\E[|T|\mid S_0=0]
 +\epsilon_0\widehat{\mathcal R}(A).
\]
The regular linear term vanishes because $\sum_A\delta_j=0$. On the
balanced slice the selected parity is $-\epsilon_0$.
Every other transversal coefficient comes only from the equator term
and the remainder. For a block-constant nonnegative atom, the integral
\eqref{eq:Fourier-integral} fixes the sign of every transversal
coefficient to be $\sigma_s=(-1)^{n_s/2+1}$. Define
\[
 \Lambda(x)=\sum_s a(n_s)\sigma_s\chi_{A_s}(x),
 \qquad Q_r=\E[\Lambda h].
\]
Apply \eqref{eq:outside-gap} to the unused blocks,
conditionally on the selected coordinates. Since
$\E|b+U|\ge\E|U|$, summing those coefficients and using
$a_{\max}/a(D_0)\le\alpha\le3/2$ gives
\begin{equation}\label{eq:cusp}
 Q_r\le\sqrt{D_0}a-\frac{a(D_0)p_0}{8}
 \E[|T|\mid S_0=0]+|\E[\Lambda\mathcal R]|.
\end{equation}
Here the equator coefficients other than the selected one were bounded
in absolute value, while the remainder was kept as a single signed
weighted sum.

For a zero-sum vector $\delta$ on a balanced slice,
\[
 \E[(\delta\cdot X)^2\mid S_0=0]=\frac{D_0}{D_0-1}\|\delta\|^2,
\]
\[
 \E[(\delta\cdot X)^4\mid S_0=0]
 =\frac{3D_0(D_0-2)}{(D_0-1)(D_0-3)}\|\delta\|^4
 -\frac{2D_0^2}{(D_0-1)(D_0-3)}\sum_j\delta_j^4
 \le12\|\delta\|^4.
\]
These identities follow by expanding and using the two- and four-point
moments of the uniform balanced slice. Adding the independent outside
sum gives conditional second moment at least $\tau^2$ and fourth moment
at most $12\tau^4$. Thus
\[
 \E[|T|\mid S_0=0]\ge\tau/\sqrt{12},\qquad
 p_0\ge1/\sqrt{2D_0}.
\]
The cusp loss in \eqref{eq:cusp} is at least
$a(D_0)\tau/(8\sqrt{24D_0})\ge\sqrt{D_0}\tau/(8\sqrt{24})$.
To compare the remainder, rename rails temporarily so that the smaller
rail in each pair comes first. Then
$n_s=n_{\min}+\sum_i\gamma_i s_i$, with even $\gamma_i\ge0$.
The Bernstein representation \eqref{eq:Bernstein}, proved independently
in Appendix~\ref{app:Bernstein}, implies that every nonconstant Walsh
coefficient of $w_s=a(n_s)$ is nonpositive. The signs $\sigma_s$ amount
to a logical translation, and the factor $b(x)$ has absolute value one.
Consequently $|\Lambda|$ is constant on logical fibers and
\[
 \|\Lambda\|_1=\sum_z|\widehat w(z)|
 =2\E_s w_s-w_0\le2a_{\max}\le6n_{\max}.
\]
Here $n_{\max}\le m(R+\Delta)\le2D_0$. Apply the uniform conditional
bound \eqref{eq:remainder} on each logical fiber to obtain
\[
 |\E[\Lambda\mathcal R]|
 \le6n_{\max}\frac{2^{14}e^{-192}\tau}{\sqrt{D_0}}
 \le2^{18}e^{-192}\sqrt{D_0}\tau
 <\frac{\sqrt{D_0}\tau}{16\sqrt{24}}.
\]
In particular, every localized balanced atom satisfies
$Q_r\le\sqrt{D_0}a$, strictly if $\tau>0$.
For a possible violation $Q_r\ge1$, take $\varepsilon=0$ in the
localization. If $\tau>0$, the cusp comparison yields
\[
 Q_r<\sqrt{D_0}a\le\sqrt{D_0a^2+\tau^2}=1,
\]
a contradiction. The case $\tau=0$ is precisely a balanced selected
parity and attains equality.

For arbitrary incoming weights, equality in block balancing forces every
independently sampled balanced vector to have this form. A partially
zero rail, two nonzero opposite rails, or unequal selected coordinate
magnitudes would produce a positive-probability strict sample. Thus the
same equality classification holds before balancing, and bias and
irrelevant weights must vanish. This proves \eqref{eq:atom-norm}; signed
parity averaging gives \eqref{eq:exact-code-norm}.

\subsection{Polynomial squared-distance stability}

Let $\mathcal A$ be the union of balanced selected-transversal parameter
rays, including zero and all input sign choices. For balanced physical
vectors $v$, the preceding localization and cusp proof gives
\begin{equation}\label{eq:balanced-stability}
 \dist(v,\mathcal A)^2\le100R\|v\|(\|v\|-Q(h_v)).
\end{equation}
Here are the constants. Normalize $\|v\|=1$ and put
$\varepsilon=1-Q(h_v)$. If $\varepsilon\ge\Delta/(100R)$, use the zero
ray, since $100R\varepsilon\ge\Delta\ge2$.
Otherwise the localization and H\"older bootstrap apply with this
$\varepsilon$, so \eqref{eq:tau-localization} and the same local cusp
comparison give $Q(h_v)\le\sqrt{D_0}a$. Consequently
$\dist(v,\mathcal A)^2\le\tau^2\le2\varepsilon$.

We now transfer this bound before block balancing without an exponentially
small sampling probability. Assume nonnegative relevant weights and
independently sample one coordinate magnitude $U_i^j$ from each rail;
let $V$ repeat this value throughout the rail. Put
$a_i=\max(U_i^0,U_i^1)$, $b_i=\min(U_i^0,U_i^1)$, and
$\bar a=m^{-1}\sum_i a_i$. The squared scalar distance to selecting one
rail of common magnitude is
\[
 d_0(U)^2=\sum_i b_i^2+\sum_i(a_i-\bar a)^2.
\]
Write $\eta=\E d_0(U)^2$ and $r_{\max}=\max r_i^j\le2R$.
For independent rail samples $X,Y$ with $A=\E X^2\ge B=\E Y^2$
and $A+B>0$ (the zero case is immediate),
\[
 \E\min(X^2,Y^2)
 \ge\frac{\E X^2Y^2}{r_{\max}(A+B)}
 \ge\frac{B}{2r_{\max}}.
\]
Choose deterministically in every pair the rail of larger second moment.
Call its sample $Z_i$ and the discarded sample $T_i$. Then
$\E\sum T_i^2\le2r_{\max}\eta$ and $|Z_i-a_i|\le T_i$.
If $P$ projects off the constant vector in $\R^m$, independence of the
$Z_i$ gives
\[
 \min_\lambda\E\sum_i(Z_i-\lambda)^2
 \le\frac{m}{m-1}\E\|PZ\|^2
 \le(4+8r_{\max})\eta.
\]
Adding discarded rails and restoring physical rail-size weights yields
\begin{equation}\label{eq:sampling-rounding}
 \dist(w,\mathcal A)^2\le24r_{\max}\E\dist(V,\mathcal A)^2.
\end{equation}

Normalize the original augmented vector to one and again put
$\varepsilon=1-Q(h_u)$. For $\varepsilon<1/2$, the block-balancing
geometric-mean inequality and $\E\|V\|^2\le1$ imply
\[
 \E\log\frac{\|V\|}{Q(h_V)}\le-\log(1-\varepsilon)\le2\varepsilon,
 \qquad
 \E\left[1-\frac{Q(h_V)}{\|V\|}\right]\le2\varepsilon.
\]
Also $\|V\|^2\le r_{\max}$ pointwise. Equations
\eqref{eq:balanced-stability} and \eqref{eq:sampling-rounding} therefore
give
\[
 \dist(w,\mathcal A)^2\le48(100R)r_{\max}^2\varepsilon.
\]
The bias and irrelevant squared norm are at most $2\varepsilon$, since
$Q(h_u)$ is bounded by the relevant norm. Altogether
\[
 \dist(u,\mathcal A)^2\le2^{16}R^3\varepsilon.
\]
For $\varepsilon\ge1/2$, the zero ray suffices. Homogeneity proves the
corresponding statement at arbitrary norm. Finally ReLU is
$1$-Lipschitz and $\E[\bar x\bar x^\top]=I$, so parameter squared
distance bounds function squared $L^2$ distance. Every face atom lies in
$\mathcal D$, proving \eqref{eq:atom-stability}.

\section{An exact Bernstein representation}\label{app:Bernstein}

The gamma-function identity for \eqref{eq:parity-cost} is
\[
 a(r)=2\sqrt{\pi r}\frac{\Gamma(r/2)}{\Gamma((r-1)/2)}
     =2\sqrt{2\pi}\,A(r/2),\qquad
 A(z)=\sqrt z\frac{\Gamma(z)}{\Gamma(z-1/2)}.
\]
Define
\[
 k(t)=\sqrt{\frac{t}{1-e^{-t}}},\qquad
 B(t)=\frac1\pi\int_0^1\frac{k(tu)}{\sqrt{u(1-u)}}\,du,
 \qquad h(t)=\tfrac12B'(t)-B''(t).
\]
The limits at zero are $k(0)=B(0)=1$ and $B'(0)=1/8$.
For $t>0$, let
\[
 D(t)=\frac1t-\frac1{e^t-1}.
\]
Then $0<D(t)<1/2$ and
$D'(t)=-1/t^2+1/[4\sinh^2(t/2)]<0$. Since
\[
 k'/k=D/2,\qquad k''/k=D^2/4+D'/2<D/8,
\]
we have $B'>0$ and $B''<B'/4$. In particular $h>0$.

Write $\mathcal L$ for the Laplace transform and
$b_0(t)=(1-e^{-t})^{-1/2}=t^{-1/2}k(t)$. The beta integral and convolution
give
\[
 \frac{\Gamma(z)}{\Gamma(z+1/2)}=\frac1{\sqrt\pi}\mathcal L[b_0](z),
 \qquad
 B=\frac1\pi(b_0*t^{-1/2}),\qquad
 z\mathcal L[B](z)=1+\mathcal L[B'](z).
\]
Therefore integration by parts yields
\begin{align*}
 A(z)&=(z-1/2)z\mathcal L[B](z)\\
 &=z-3/8+\mathcal L[B''-B'/2](z)
 =z-3/8-\int_0^\infty e^{-zt}h(t)\,dt.
\end{align*}
All functions and derivatives have at most polynomial growth, so these
integrals converge for $z>0$. At $z=1/2$, $A$ extends continuously to zero.
Consequently
\[
 A(x+1/2)=x+\int_0^\infty e^{-t/2}(1-e^{-xt})h(t)\,dt
 \quad(x\ge0),
\]
which is the claimed strict Bernstein representation.

For the crude bound $h\le3$, use
$k(t)\le\sqrt{1+t}$ and $D(t)\le\min(1/2,1/t)$. For $0<t\le1$,
\[
 |D'(t)|=\frac{2(\cosh t-1)-t^2}{4t^2\sinh^2(t/2)}
 \le2(\cosh1-1)-1<1;
\]
the numerator's power series starts at degree four and the denominator
is at least $t^4$. For $t\ge1$, $|D'(t)|\le t^{-2}$.
The displayed derivative identities then give $k'\le1$ and
$|k''|\le2$ on the whole positive line. Hence $B'\le1$,
$|B''|\le2$, and $0<h\le5/2<3$. This proves
\eqref{eq:Bernstein} with all bounds used above.

\end{document}
